% !TeX program = lualatex
\documentclass[12pt,aspectratio=169]{beamer}
\usepackage[utf8]{inputenc}
\usepackage[T1]{fontenc}
\usepackage{amsmath,amsfonts,amssymb}
\usepackage{graphicx}
\usepackage{tikz}
\usepackage{xcolor}
\usepackage{hyperref}
\usepackage{anyfontsize}
\usepackage{lmodern}

% ====== MAKE ALL FONTS SMALLER ======
\renewcommand{\tiny}{\fontsize{5}{6}\selectfont}
\renewcommand{\scriptsize}{\fontsize{6}{7}\selectfont}
\renewcommand{\footnotesize}{\fontsize{7}{8}\selectfont}
\renewcommand{\small}{\fontsize{8}{9}\selectfont}
\renewcommand{\normalsize}{\fontsize{9}{10}\selectfont}
\renewcommand{\large}{\fontsize{10}{11}\selectfont}
\renewcommand{\Large}{\fontsize{11}{13}\selectfont}
\renewcommand{\LARGE}{\fontsize{13}{15}\selectfont}
\renewcommand{\huge}{\fontsize{15}{17}\selectfont}
\renewcommand{\Huge}{\fontsize{18}{21}\selectfont}

% ====== COLOR THEME ======
\definecolor{redcorrect}{RGB}{200,30,30}
\definecolor{bluecorrect}{RGB}{30,80,200}
\definecolor{greencheck}{RGB}{0,150,50}
\definecolor{lightgray}{RGB}{245,245,245}
\definecolor{darkgray}{RGB}{50,50,50}
\definecolor{softyellow}{RGB}{255,248,200}
\definecolor{infoblue}{RGB}{230,240,252}
\definecolor{steppurple}{RGB}{150,50,200}
\definecolor{deepblue}{RGB}{20,60,150}
\definecolor{orangewarning}{RGB}{255,140,0}

\setbeamercolor{title}{fg=white,bg=redcorrect}
\setbeamercolor{frametitle}{fg=white,bg=redcorrect}
\setbeamercolor{block title}{fg=white,bg=bluecorrect}
\setbeamercolor{block body}{fg=darkgray,bg=lightgray}
\setbeamercolor{normal text}{fg=darkgray}
\usecolortheme{default}

% ====== CUSTOM COMMANDS ======
\newcommand{\correct}[1]{\textcolor{greencheck}{\textbf{\checkmark}} \textcolor{greencheck}{#1}}
\newcommand{\highlightred}[1]{\textcolor{redcorrect}{\textbf{#1}}}
\newcommand{\highlightblue}[1]{\textcolor{bluecorrect}{\textbf{#1}}}
\newcommand{\highlightorange}[1]{\textcolor{orangewarning}{\textbf{#1}}}
\newcommand{\tipbox}[1]{\fcolorbox{bluecorrect}{softyellow}{\parbox{0.88\textwidth}{\centering #1}}}
\newcommand{\steptitle}[1]{\textcolor{steppurple}{\textbf{#1}}}
\newcommand{\keypoint}[1]{\textcolor{deepblue}{\textbf{#1}}}
\newcommand{\redtitle}[1]{\textcolor{redcorrect}{\textbf{#1}}}

\newcommand{\stepbox}[1]{%
	\begin{center}
		\fcolorbox{darkgray}{white}{%
			\parbox{0.90\textwidth}{\vspace{0.1cm} #1 \vspace{0.1cm}}%
		}
	\end{center}
}

\begin{document}

% ================================================================
% SLIDE 1: TITLE
% ================================================================
\begin{frame}
\centering
\vspace{0.8cm}
{\fontsize{20}{24}\selectfont \textbf{Money Laundering and Financial Crime}}\\[0.2cm]
{\fontsize{15}{18}\selectfont Complete Tutorial Guide — All Topics}\\[0.4cm]
{\fontsize{12}{14}\selectfont Step-by-Step Learning for Every Subject}\\[0.4cm]
{\fontsize{10}{12}\selectfont Compliance Training Department}
\vspace{0.3cm}
\rule{5cm}{1pt}
\end{frame}

% ================================================================
% SLIDE 2: TUTORIAL 1 — CHINA AML REGULATIONS
% ================================================================
\begin{frame}
\frametitle{\fontsize{14}{17}\selectfont TUTORIAL 1: Mainland China AML Regulations}

\small
\begin{block}{\fontsize{11}{13}\selectfont Topic: AFC Regulations and Regimes}
\footnotesize
Mainland China's AML regulations have specific features that are important for compliance professionals to understand.
\end{block}

\vspace{0.03cm}

\stepbox{\footnotesize
\steptitle{Step 1:} Understand the scope of China's AML regulations. They have evolved to address modern financial crime challenges. \\
\steptitle{Step 2:} Learn the key features. Two important features are: \\
\quad (1) \textbf{Provisions restricting excessive risk measures} \\
\quad (2) \textbf{Extraterritoriality application} \\
\steptitle{Step 3:} Understand why these features matter. \\
\quad - Excessive risk measures: Prevents overreach that could harm legitimate business \\
\quad - Extraterritoriality: China can enforce AML rules on activities outside its borders \\
\steptitle{Step 4:} See what is NOT a feature. \textbf{Exemption for cryptocurrency activities} is incorrect — China has strict crypto regulations. \\
\steptitle{Step 5:} Apply the lesson. China's AML regime includes restrictions on excessive measures and extraterritorial application.}

\vspace{0.08cm}
\tipbox{\footnotesize \textbf{CORE LESSON:} China's AML regulations include provisions restricting excessive risk measures and extraterritoriality application.}
\end{frame}

% ================================================================
% SLIDE 3: TUTORIAL 2 — UK TRADE-BASED ML FAILURE
% ================================================================
\begin{frame}
\frametitle{\fontsize{14}{17}\selectfont TUTORIAL 2: Trade-Based Money Laundering Failures}

\small
\begin{block}{\fontsize{11}{13}\selectfont Topic: Trade Finance Compliance}
\footnotesize
Financial institutions can receive fines for inadequate transaction monitoring. Understanding the specific weaknesses is essential.
\end{block}

\vspace{0.03cm}

\stepbox{\footnotesize
\steptitle{Step 1:} Understand the scenario. A UK financial institution receives a regulatory fine for inadequate transaction monitoring tied to trade-based money laundering. \\
\steptitle{Step 2:} Recognize the weakness that contributed. The weakness was \textbf{failure to recognize high-risk typologies within client sectors}. \\
\steptitle{Step 3:} Understand why this is a weakness. \\
\quad - Different sectors have different TBML typologies \\
\quad - Failing to recognize them means missing suspicious activity \\
\quad - Regulators expect sector-specific knowledge \\
\steptitle{Step 4:} Apply the lesson. Institutions must understand and monitor for TBML typologies specific to their client sectors.}

\vspace{0.08cm}
\tipbox{\footnotesize \textbf{CORE LESSON:} TBML fines often result from failure to recognize high-risk typologies within specific client sectors.}
\end{frame}

% ================================================================
% SLIDE 4: TUTORIAL 3 — INCLUSIVE BANKING PRACTICES
% ================================================================
\begin{frame}
\frametitle{\fontsize{14}{17}\selectfont TUTORIAL 3: Inclusive Banking Practices}

\small
\begin{block}{\fontsize{11}{13}\selectfont Topic: Financial Inclusion}
\footnotesize
Inclusive banking strategies help reach underserved populations while maintaining compliance.
\end{block}

\vspace{0.03cm}

\stepbox{\footnotesize
\steptitle{Step 1:} Understand the goal of inclusive banking. It aims to bring more people into the formal financial system. \\
\steptitle{Step 2:} Learn the practices that support inclusion. Two key practices are: \\
\quad (1) \textbf{Offering multilingual customer support options} \\
\quad (2) \textbf{Developing mobile apps with offline functionality} \\
\steptitle{Step 3:} Understand why these support inclusion. \\
\quad - Multilingual support: Reaches diverse populations \\
\quad - Offline functionality: Reaches areas with poor connectivity \\
\steptitle{Step 4:} Apply the lesson. Inclusive banking removes barriers to access through language and technology.}

\vspace{0.08cm}
\tipbox{\footnotesize \textbf{CORE LESSON:} Inclusive banking practices include multilingual support and mobile apps with offline functionality.}
\end{frame}

% ================================================================
% SLIDE 5: TUTORIAL 4 — DORA TESTING REQUIREMENTS
% ================================================================
\begin{frame}
\frametitle{\fontsize{14}{17}\selectfont TUTORIAL 4: Digital Operational Resilience Act}

\small
\begin{block}{\fontsize{11}{13}\selectfont Topic: DORA Requirements}
\footnotesize
The Digital Operational Resilience Act (DORA) requires financial entities to test their ICT systems.
\end{block}

\vspace{0.03cm}

\stepbox{\footnotesize
\steptitle{Step 1:} Understand what DORA is. It is an EU regulation focused on ICT resilience in financial services. \\
\steptitle{Step 2:} Learn the testing requirement. DORA requires advanced testing \textbf{to detect potential vulnerabilities and assess system defenses before incidents occur}. \\
\steptitle{Step 3:} Understand why this matters. \\
\quad - Proactive testing finds weaknesses before attackers do \\
\quad - Assessment of defenses ensures they are effective \\
\quad - Prevention is better than response \\
\steptitle{Step 4:} Apply the lesson. DORA emphasizes proactive vulnerability detection through advanced testing.}

\vspace{0.08cm}
\tipbox{\footnotesize \textbf{CORE LESSON:} DORA requires advanced ICT testing to detect vulnerabilities and assess defenses before incidents occur.}
\end{frame}

% ================================================================
% SLIDE 6: TUTORIAL 5 — UK BRIBERY ACT vs FCPA
% ================================================================
\begin{frame}
\frametitle{\fontsize{14}{17}\selectfont TUTORIAL 5: UK Bribery Act vs US FCPA}

\small
\begin{block}{\fontsize{11}{13}\selectfont Topic: Anti-Corruption Laws}
\footnotesize
The UK Bribery Act 2010 and the US Foreign Corrupt Practices Act (FCPA) are key anti-corruption laws with important differences.
\end{block}

\vspace{0.03cm}

\stepbox{\footnotesize
\steptitle{Step 1:} Understand the UK Bribery Act. It is considered one of the toughest anti-bribery laws in the world. \\
\steptitle{Step 2:} Learn the key distinction. \textbf{The UK Bribery Act 2010 prohibits facilitation payments made to officials}. \\
\steptitle{Step 3:} Compare to the FCPA. The FCPA has a defense for facilitation payments (small payments to speed routine government actions). The UK Bribery Act does not. \\
\steptitle{Step 4:} Understand the implication. \\
\quad - UK Bribery Act: No facilitation payments allowed \\
\quad - FCPA: Facilitation payments may be allowed \\
\steptitle{Step 5:} Apply the lesson. The UK Bribery Act is stricter than the FCPA — it prohibits all facilitation payments.}

\vspace{0.08cm}
\tipbox{\footnotesize \textbf{CORE LESSON:} The UK Bribery Act prohibits facilitation payments, making it stricter than the FCPA.}
\end{frame}

% ================================================================
% SLIDE 7: TUTORIAL 6 — BATCH SCREENING RESPONSE
% ================================================================
\begin{frame}
\frametitle{\fontsize{14}{17}\selectfont TUTORIAL 6: Batch Screening and Risk Changes}

\small
\begin{block}{\fontsize{11}{13}\selectfont Topic: Sanctions Screening}
\footnotesize
When a customer is flagged in a sanctions list and adverse media, the risk profile has changed and action is required.
\end{block}

\vspace{0.03cm}

\stepbox{\footnotesize
\steptitle{Step 1:} Understand the scenario. A customer previously categorized as low risk is now simultaneously flagged in a recent sanctions list and associated with adverse media coverage. \\
\steptitle{Step 2:} Recognize what has happened. The customer's risk profile has \textbf{fundamentally changed}. Two red flags have appeared. \\
\steptitle{Step 3:} Understand the required response. The institution must: \\
\quad - Investigate the sanctions alert \\
\quad - Review the adverse media \\
\quad - Reassess the customer risk \\
\quad - Consider filing a SAR if suspicious \\
\steptitle{Step 4:} Apply the lesson. When a customer is flagged in sanctions and adverse media, immediate investigation and reassessment are required.}

\vspace{0.08cm}
\tipbox{\footnotesize \textbf{CORE LESSON:} When a customer is flagged in sanctions lists and adverse media, the risk profile has changed and immediate investigation is required.}
\end{frame}

% ================================================================
% SLIDE 8: TUTORIAL 7 — EU AML REGULATORY CHALLENGES
% ================================================================
\begin{frame}
\frametitle{\fontsize{14}{17}\selectfont TUTORIAL 7: EU AML Regulatory Challenges}

\small
\begin{block}{\fontsize{11}{13}\selectfont Topic: EU AML Framework}
\footnotesize
The EU AML framework faces specific challenges in implementation across member states.
\end{block}

\vspace{0.03cm}

\stepbox{\footnotesize
\steptitle{Step 1:} Understand the EU framework. EU AML directives set minimum standards that member states must implement. \\
\steptitle{Step 2:} Recognize the challenges. Three key challenges are: \\
\quad (1) \textbf{Fragmented supervision of cross-border financial activity} \\
\quad (2) \textbf{Lack of cooperation among regulatory authorities} \\
\quad (3) \textbf{Gaps in transposition of directives into national law} \\
\steptitle{Step 3:} Understand why these are challenges. \\
\quad - Fragmented supervision: No single oversight body \\
\quad - Lack of cooperation: Information doesn't flow across borders \\
\quad - Transposition gaps: Different countries implement differently \\
\steptitle{Step 4:} Apply the lesson. The EU AML framework is weakened by fragmentation, lack of cooperation, and implementation gaps.}

\vspace{0.08cm}
\tipbox{\footnotesize \textbf{CORE LESSON:} EU AML challenges include fragmented supervision, lack of regulatory cooperation, and gaps in directive transposition.}
\end{frame}

% ================================================================
% SLIDE 9: TUTORIAL 8 — AMLR AND 6AMLD EXTENSIONS
% ================================================================
\begin{frame}
\frametitle{\fontsize{14}{17}\selectfont TUTORIAL 8: AMLR and 6AMLD Extensions}

\small
\begin{block}{\fontsize{11}{13}\selectfont Topic: EU AML Directives}
\footnotesize
The AMLR and 6AMLD extended AML obligations to new types of entities and professionals.
\end{block}

\vspace{0.03cm}

\stepbox{\footnotesize
\steptitle{Step 1:} Understand what AMLR and 6AMLD are. They are EU AML regulations that expanded the scope of AML obligations. \\
\steptitle{Step 2:} Learn the entities covered. Two key extensions are: \\
\quad (1) \textbf{Cryptoasset service providers} \\
\quad (2) \textbf{Soccer agents} \\
\steptitle{Step 3:} Understand why these were added. \\
\quad - Cryptoasset providers: Emerging risk area \\
\quad - Soccer agents: High-value transfers, potential for laundering \\
\steptitle{Step 4:} Apply the lesson. AML obligations now extend to cryptoasset providers and soccer agents.}

\vspace{0.08cm}
\tipbox{\footnotesize \textbf{CORE LESSON:} AMLR and 6AMLD extended AML obligations to cryptoasset service providers and soccer agents.}
\end{frame}

% ================================================================
% SLIDE 10: TUTORIAL 9 — FOREIGN EMBASSY ACCOUNT RISK
% ================================================================
\begin{frame}
\frametitle{\fontsize{14}{17}\selectfont TUTORIAL 9: Foreign Embassy Account Risk Assessment}

\small
\begin{block}{\fontsize{11}{13}\selectfont Topic: DNFBP — Embassy Accounts}
\footnotesize
Accounts held by foreign embassies and missions present specific AML risks that require ongoing monitoring.
\end{block}

\vspace{0.03cm}

\stepbox{\footnotesize
\steptitle{Step 1:} Understand the risk. Embassy accounts can be used for money laundering or terrorist financing. \\
\steptitle{Step 2:} Learn the actions for ongoing risk evaluation. Three key actions are: \\
\quad (1) \textbf{Monitor transaction volume and compare it with the entity's stated purpose} \\
\quad (2) \textbf{Review account activity periodically for evidence of third-party or private use} \\
\quad (3) \textbf{Require clarification if funds appear to originate from or are sent to sanctioned jurisdictions} \\
\steptitle{Step 3:} Understand why these are important. \\
\quad - Volume monitoring: Identifies unusual activity \\
\quad - Third-party review: Prevents misuse \\
\quad - Sanction checks: Ensures compliance \\
\steptitle{Step 4:} Apply the lesson. Embassy accounts require ongoing monitoring of volume, third-party use, and sanctions exposure.}

\vspace{0.08cm}
\tipbox{\footnotesize \textbf{CORE LESSON:} Evaluate embassy account risk through volume monitoring, third-party use review, and sanctions jurisdiction checks.}
\end{frame}

% ================================================================
% SLIDE 11: TUTORIAL 10 — TRADE-BASED MISCLASSIFICATION
% ================================================================
\begin{frame}
\frametitle{\fontsize{14}{17}\selectfont TUTORIAL 10: Trade-Based Misclassification}

\small
\begin{block}{\fontsize{11}{13}\selectfont Topic: Trade-Based Money Laundering}
\footnotesize
Misclassification of goods in Free-Trade Zones is a common money laundering tactic.
\end{block}

\vspace{0.03cm}

\stepbox{\footnotesize
\steptitle{Step 1:} Understand the scenario. A logistics firm operating out of an FTZ lists pharmaceuticals as incoming cargo but later categorizes outbound goods as shoes. \\
\steptitle{Step 2:} Recognize the tactic. This is \textbf{trade-based misclassification masking value movement}. \\
\steptitle{Step 3:} Understand how it works. \\
\quad - Incoming: Pharmaceuticals (high value) \\
\quad - Outbound: Shoes (lower value) \\
\quad - The difference in value is moved illicitly \\
\steptitle{Step 4:} See why this is effective. \\
\quad - Changes product classification \\
\quad - Changes value \\
\quad - Moves money without detection \\
\steptitle{Step 5:} Apply the lesson. Goods misclassification in FTZs is a red flag for trade-based laundering.}

\vspace{0.08cm}
\tipbox{\footnotesize \textbf{CORE LESSON:} Trade-based misclassification — changing product categorization — masks value movement and facilitates laundering.}
\end{frame}

% ================================================================
% SLIDE 12: TUTORIAL 11 — BLOCKCHAIN ANALYTICS FEATURES
% ================================================================
\begin{frame}
\frametitle{\fontsize{14}{17}\selectfont TUTORIAL 11: Effective Blockchain Analytics}

\small
\begin{block}{\fontsize{11}{13}\selectfont Topic: Cryptoasset Compliance}
\footnotesize
Effective blockchain analytics for compliance requires specific features to identify and monitor risk.
\end{block}

\vspace{0.03cm}

\stepbox{\footnotesize
\steptitle{Step 1:} Understand the need for blockchain analytics. Cryptoassets require specialized monitoring. \\
\steptitle{Step 2:} Learn the characteristics of effective analytics. Three key characteristics are: \\
\quad (1) \textbf{Coverage across all major blockchains} \\
\quad (2) \textbf{Integration with robust entity resolution tools} \\
\quad (3) \textbf{Broad alerting functionality for potential risk behaviors} \\
\steptitle{Step 3:} Understand why these matter. \\
\quad - Coverage: All blockchain data \\
\quad - Entity resolution: Connects addresses to identities \\
\quad - Alerting: Detects suspicious behavior \\
\steptitle{Step 4:} Apply the lesson. Effective blockchain analytics requires comprehensive coverage, entity resolution, and alerting.}

\vspace{0.08cm}
\tipbox{\footnotesize \textbf{CORE LESSON:} Effective blockchain analytics requires coverage across blockchains, entity resolution integration, and broad alerting.}
\end{frame}

% ================================================================
% SLIDE 13: TUTORIAL 12 — PAYMENT SERVICE PROVIDER RISKS
% ================================================================
\begin{frame}
\frametitle{\fontsize{14}{17}\selectfont TUTORIAL 12: Payment Service Provider Risks}

\small
\begin{block}{\fontsize{11}{13}\selectfont Topic: Payment Service Providers}
\footnotesize
Payment service providers face specific AML risks that must be managed.
\end{block}

\vspace{0.03cm}

\stepbox{\footnotesize
\steptitle{Step 1:} Understand the PSP risk environment. PSPs process large volumes of transactions for many customers. \\
\steptitle{Step 2:} Learn the relevant risks. Two key risks are: \\
\quad (1) \textbf{Risk from operating in jurisdictions with limited regulatory oversight} \\
\quad (2) \textbf{Risk of facilitating fraudulent transactions} \\
\steptitle{Step 3:} Understand why these are relevant. \\
\quad - Limited oversight: Weak AML controls \\
\quad - Fraud facilitation: PSPs can be unwitting conduits \\
\steptitle{Step 4:} Apply the lesson. PSPs must manage jurisdictional risk and fraud facilitation risk.}

\vspace{0.08cm}
\tipbox{\footnotesize \textbf{CORE LESSON:} Payment service providers must manage risks from operating in weak oversight jurisdictions and facilitating fraud.}
\end{frame}

% ================================================================
% SLIDE 14: TUTORIAL 13 — NFT LAYERING
% ================================================================
\begin{frame}
\frametitle{\fontsize{14}{17}\selectfont TUTORIAL 13: NFT Layering}

\small
\begin{block}{\fontsize{11}{13}\selectfont Topic: Digital Assets and Laundering}
\footnotesize
NFTs can be used for money laundering layering through inflated and inconsistent pricing.
\end{block}

\vspace{0.03cm}

\stepbox{\footnotesize
\steptitle{Step 1:} Understand the scenario. A user repeatedly buys digital art NFTs with stablecoin and resells them at inflated and inconsistent prices. \\
\steptitle{Step 2:} Recognize the suspicion. The user is \textbf{using NFTs to layer illicit funds}. \\
\steptitle{Step 3:} Understand how it works. \\
\quad - Buy NFT with illicit funds \\
\quad - Sell NFT at inflated price \\
\quad - Receive "clean" proceeds \\
\steptitle{Step 4:} See why this is layering. \\
\quad - NFT sale appears legitimate \\
\quad - Price manipulation hides the source \\
\quad - Funds appear clean \\
\steptitle{Step 5:} Apply the lesson. Inflated and inconsistent NFT pricing is a red flag for layering.}

\vspace{0.08cm}
\tipbox{\footnotesize \textbf{CORE LESSON:} Inflated NFT prices and inconsistent pricing suggest layering of illicit funds through digital art transactions.}
\end{frame}

% ================================================================
% SLIDE 15: TUTORIAL 14 — MSB AGENT NETWORK RISK
% ================================================================
\begin{frame}
\frametitle{\fontsize{14}{17}\selectfont TUTORIAL 14: MSB Agent Network Risk}

\small
\begin{block}{\fontsize{11}{13}\selectfont Topic: Money Service Businesses}
\footnotesize
An MSB's agent network influences its operational and geographic risk exposure.
\end{block}

\vspace{0.03cm}

\stepbox{\footnotesize
\steptitle{Step 1:} Understand the scenario. An onboarding questionnaire asks about an MSB's relationships with third-party agents. \\
\steptitle{Step 2:} Recognize why this is critical. \textbf{The MSB's agent network could influence both operational and geographic risk exposure}. \\
\steptitle{Step 3:} Understand the risk. \\
\quad - Agents may have weak controls \\
\quad - Agents may be in high-risk locations \\
\quad - The MSB is responsible for agents \\
\steptitle{Step 4:} Apply the lesson. MSB agent networks must be assessed because they create operational and geographic risk.}

\vspace{0.08cm}
\tipbox{\footnotesize \textbf{CORE LESSON:} MSB agent networks influence operational and geographic risk exposure and must be assessed.}
\end{frame}

% ================================================================
% SLIDE 16: TUTORIAL 15 — OVER/UNDER-INVOICING
% ================================================================
\begin{frame}
\frametitle{\fontsize{14}{17}\selectfont TUTORIAL 15: Over/Under-Invoicing}

\small
\begin{block}{\fontsize{11}{13}\selectfont Topic: Trade-Based Laundering}
\footnotesize
Vastly different pricing for identical cargo is a strong indicator of over/under-invoicing.
\end{block}

\vspace{0.03cm}

\stepbox{\footnotesize
\steptitle{Step 1:} Understand the scenario. A freight forwarder reports shipping clients declaring identical cargo with identical weight but with vastly different pricing across invoices. \\
\steptitle{Step 2:} Recognize the risk. This is \textbf{over/under-invoicing to obscure monetary transfers}. \\
\steptitle{Step 3:} Understand how it works. \\
\quad - Same goods, different prices \\
\quad - Over-invoice: Move money out \\
\quad - Under-invoice: Move money in \\
\steptitle{Step 4:} Apply the lesson. Different prices for identical goods is a red flag for trade-based money laundering.}

\vspace{0.08cm}
\tipbox{\footnotesize \textbf{CORE LESSON:} Vastly different pricing for identical cargo indicates over/under-invoicing to obscure monetary transfers.}
\end{frame}

% ================================================================
% SLIDE 17: TUTORIAL 16 — LAW FIRM UBO DISCLOSURE
% ================================================================
\begin{frame}
\frametitle{\fontsize{14}{17}\selectfont TUTORIAL 16: Law Firm UBO Disclosure}

\small
\begin{block}{\fontsize{11}{13}\selectfont Topic: Legal Sector CDD}
\footnotesize
Law firms must decline services when clients refuse to disclose beneficial ownership information.
\end{block}

\vspace{0.03cm}

\stepbox{\footnotesize
\steptitle{Step 1:} Understand the scenario. A client asks a law firm to incorporate a company in a low-transparency country and requests nominee shareholders. The client refuses to disclose beneficiary identity. \\
\steptitle{Step 2:} Recognize the red flags. \\
\quad - Low-transparency jurisdiction \\
\quad - Nominee shareholders \\
\quad - Refusal to disclose UBO \\
\steptitle{Step 3:} Understand the correct action. The firm should \textbf{decline the service due to insufficient CDD and refusal to disclose UBO}. \\
\steptitle{Step 4:} Apply the lesson. Law firms must decline services when clients refuse to provide beneficial ownership information.}

\vspace{0.08cm}
\tipbox{\footnotesize \textbf{CORE LESSON:} Law firms must decline services when clients refuse to disclose beneficial ownership — CDD is mandatory.}
\end{frame}

% ================================================================
% SLIDE 18: TUTORIAL 17 — DUAL-USE EXPORT CONTROLS
% ================================================================
\begin{frame}
\frametitle{\fontsize{14}{17}\selectfont TUTORIAL 17: Dual-Use Export Controls}

\small
\begin{block}{\fontsize{11}{13}\selectfont Topic: Export Controls and Sanctions}
\footnotesize
Dual-use goods require careful screening, especially when destined for sanctioned countries.
\end{block}

\vspace{0.03cm}

\stepbox{\footnotesize
\steptitle{Step 1:} Understand the scenario. A shipper is contracted to deliver radar components to a foreign government's defense ministry. The goods are dual-use and the destination country is under partial EU sanctions. \\
\steptitle{Step 2:} Recognize the compliance requirement. The shipper should \textbf{suspend delivery and initiate an export control and sanctions screening}. \\
\steptitle{Step 3:} Understand why this is required. \\
\quad - Dual-use goods: Can be used for military purposes \\
\quad - Sanctions: Restrictions on the destination \\
\steptitle{Step 4:} Apply the lesson. Dual-use goods to sanctioned destinations require export control and sanctions screening.}

\vspace{0.08cm}
\tipbox{\footnotesize \textbf{CORE LESSON:} Dual-use goods to sanctioned destinations require suspending delivery and initiating export control and sanctions screening.}
\end{frame}

% ================================================================
% SLIDE 19: TUTORIAL 18 — LAW FIRM STR FILING
% ================================================================
\begin{frame}
\frametitle{\fontsize{14}{17}\selectfont TUTORIAL 18: Law Firm Suspicious Transaction Reporting}

\small
\begin{block}{\fontsize{11}{13}\selectfont Topic: Legal Sector Reporting}
\footnotesize
Law firms must file suspicious transaction reports when transactions appear unusual.
\end{block}

\vspace{0.03cm}

\stepbox{\footnotesize
\steptitle{Step 1:} Understand the scenario. A Canadian law firm delivers real estate closing services to foreign clients. One client insists proceeds should be paid to a foreign third party unrelated to the transaction. \\
\steptitle{Step 2:} Recognize the red flag. Payment to an unrelated third party is unusual and suspicious. \\
\steptitle{Step 3:} Understand the correct action. The firm should \textbf{file a Suspicious Transaction Report and delay the transaction}. \\
\steptitle{Step 4:} Apply the lesson. Law firms must file STRs and delay transactions when third-party payments are requested without justification.}

\vspace{0.08cm}
\tipbox{\footnotesize \textbf{CORE LESSON:} Law firms must file STRs and delay transactions when clients request payment to unrelated third parties.}
\end{frame}

% ================================================================
% SLIDE 20: TUTORIAL 19 — FIU GUIDANCE AND EWRA
% ================================================================
\begin{frame}
\frametitle{\fontsize{14}{17}\selectfont TUTORIAL 19: Responding to FIU Guidance}

\small
\begin{block}{\fontsize{11}{13}\selectfont Topic: Guidance Implementation}
\footnotesize
When FIUs issue guidance on emerging risks, institutions should assess their exposure and update risk assessments.
\end{block}

\vspace{0.03cm}

\stepbox{\footnotesize
\steptitle{Step 1:} Understand the scenario. An FIU policy paper identifies that a regional real estate product has become associated with fraud networks. \\
\steptitle{Step 2:} Recognize the correct response. The institution should \textbf{assess internal risk exposure and update the EWRA accordingly}. \\
\steptitle{Step 3:} Understand why immediate suspension is not correct. \\
\quad - Immediate suspension is reactive and may not be necessary \\
\quad - Assessment first: Understand the exposure \\
\quad - Then decide on controls \\
\steptitle{Step 4:} Apply the lesson. Respond to FIU guidance by assessing risk and updating the enterprise-wide risk assessment.}

\vspace{0.08cm}
\tipbox{\footnotesize \textbf{CORE LESSON:} Respond to FIU guidance by assessing internal risk exposure and updating the EWRA before implementing controls.}
\end{frame}

% ================================================================
% SLIDE 21: TUTORIAL 20 — GUIDANCE DOCUMENT IMPACTS
% ================================================================
\begin{frame}
\frametitle{\fontsize{14}{17}\selectfont TUTORIAL 20: Guidance Document Implementation Impacts}

\small
\begin{block}{\fontsize{11}{13}\selectfont Topic: Guidance Implementation}
\footnotesize
Implementing guidance documents has specific impacts that organizations must consider.
\end{block}

\vspace{0.03cm}

\stepbox{\footnotesize
\steptitle{Step 1:} Understand the importance of considering impacts. Guidance implementation affects multiple areas of the organization. \\
\steptitle{Step 2:} Learn the important impacts to consider. Three key impacts are: \\
\quad (1) \textbf{Customer experience degradation} \\
\quad (2) \textbf{Control redundancy overlaps} \\
\quad (3) \textbf{Resource impacts related to IT and compliance functions} \\
\steptitle{Step 3:} Understand why these matter. \\
\quad - Customer experience: Can affect business \\
\quad - Control redundancy: Creates inefficiency \\
\quad - Resource impacts: Budget and staffing \\
\steptitle{Step 4:} Apply the lesson. Consider customer experience, control overlaps, and resource impacts when implementing guidance.}

\vspace{0.08cm}
\tipbox{\footnotesize \textbf{CORE LESSON:} Guidance implementation impacts include customer experience degradation, control redundancy overlaps, and resource impacts.}
\end{frame}

% ================================================================
% SLIDE 22: TUTORIAL 21 — AFC GUIDANCE DOCUMENT EFFECTIVENESS
% ================================================================
\begin{frame}
\frametitle{\fontsize{14}{17}\selectfont TUTORIAL 21: Effective Use of AFC Guidance}

\small
\begin{block}{\fontsize{11}{13}\selectfont Topic: Guidance Effectiveness}
\footnotesize
Effective use of AFC guidance documents strengthens compliance programs.
\end{block}

\vspace{0.03cm}

\stepbox{\footnotesize
\steptitle{Step 1:} Understand the purpose of AFC guidance documents. They provide direction on implementing AML/CFT controls. \\
\steptitle{Step 2:} Recognize outcomes of effective use. Effective use results in: \\
\quad - Consistent application of controls \\
\quad - Alignment with regulatory expectations \\
\quad - Demonstrated compliance \\
\steptitle{Step 3:} Understand why effectiveness matters. \\
\quad - Regulators expect guidance to be followed \\
\quad - Effective implementation reduces risk \\
\quad - It demonstrates good governance \\
\steptitle{Step 4:} Apply the lesson. Effective guidance implementation results in consistent, aligned, and demonstrable compliance.}

\vspace{0.08cm}
\tipbox{\footnotesize \textbf{CORE LESSON:} Effective use of AFC guidance documents results in consistent, aligned, and demonstrable compliance.}
\end{frame}

% ================================================================
% SLIDE 23: TUTORIAL 22 — LEGAL SECTOR RISK MITIGATION
% ================================================================
\begin{frame}
\frametitle{\fontsize{14}{17}\selectfont TUTORIAL 22: Legal Sector Risk Mitigation}

\small
\begin{block}{\fontsize{11}{13}\selectfont Topic: Legal Sector AML}
\footnotesize
The legal sector faces specific AML risks that require targeted mitigation approaches.
\end{block}

\vspace{0.03cm}

\stepbox{\footnotesize
\steptitle{Step 1:} Understand the risks in the legal sector. Lawyers handle client money and facilitate transactions. \\
\steptitle{Step 2:} Learn the effective mitigation approaches. Two key approaches are: \\
\quad (1) \textbf{Maintaining strong internal training programs on typologies and reporting} \\
\quad (2) \textbf{Restricting involvement in escrow arrangements unless risk-assessed} \\
\steptitle{Step 3:} Understand why these are effective. \\
\quad - Training: Builds awareness and skills \\
\quad - Escrow restrictions: Prevents high-risk activity \\
\steptitle{Step 4:} Apply the lesson. Legal sector risk mitigation requires training and risk assessment of escrow services.}

\vspace{0.08cm}
\tipbox{\footnotesize \textbf{CORE LESSON:} Legal sector risk mitigation includes strong training programs and risk-assessed escrow involvement.}
\end{frame}

% ================================================================
% SLIDE 24: TUTORIAL 23 — PEP USING INDIRECT OWNERSHIP
% ================================================================
\begin{frame}
\frametitle{\fontsize{14}{17}\selectfont TUTORIAL 23: PEP Using Indirect Ownership}

\small
\begin{block}{\fontsize{11}{13}\selectfont Topic: Politically Exposed Persons}
\footnotesize
PEPs using indirect ownership intermediaries present significant AML concerns.
\end{block}

\vspace{0.03cm}

\stepbox{\footnotesize
\steptitle{Step 1:} Understand the concern. A buyer is a PEP using indirect ownership intermediaries. \\
\steptitle{Step 2:} Recognize why this is concerning. \\
\quad - PEPs are high-risk \\
\quad - Indirect ownership hides the connection \\
\quad - Intermediaries complicate due diligence \\
\steptitle{Step 3:} Understand the AML perspective. \\
\quad - PEPs require enhanced due diligence \\
\quad - Indirect ownership makes it harder \\
\quad - Risk of corruption is higher \\
\steptitle{Step 4:} Apply the lesson. PEPs using indirect ownership intermediaries require enhanced scrutiny.}

\vspace{0.08cm}
\tipbox{\footnotesize \textbf{CORE LESSON:} PEPs using indirect ownership intermediaries present a significant AML concern requiring enhanced due diligence.}
\end{frame}

% ================================================================
% SLIDE 25: SUMMARY
% ================================================================
\begin{frame}
\frametitle{\fontsize{14}{17}\selectfont SUMMARY: Key Lessons from All Topics}

\small
\begin{block}{\fontsize{11}{13}\selectfont What We Have Learned:}
\footnotesize
\begin{enumerate}
    \item \textbf{China AML:} Restrictions on excessive risk measures, extraterritoriality.
    \item \textbf{UK TBML Failure:} Failure to recognize high-risk typologies in sectors.
    \item \textbf{Inclusive Banking:} Multilingual support, offline mobile apps.
    \item \textbf{DORA:} Advanced ICT testing to detect vulnerabilities.
    \item \textbf{UK Bribery Act:} Prohibits facilitation payments.
    \item \textbf{Batch Screening:} Sanctions + adverse media = immediate investigation.
    \item \textbf{EU Challenges:} Fragmented supervision, lack of cooperation, transposition gaps.
    \item \textbf{AMLR/6AMLD:} Extended to cryptoasset providers and soccer agents.
    \item \textbf{Embassy Accounts:} Monitor volume, third-party use, sanctions jurisdictions.
    \item \textbf{Trade Misclassification:} Goods misclassification masks value movement.
    \item \textbf{Blockchain Analytics:} Coverage, entity resolution, alerting.
    \item \textbf{PSP Risks:} Limited oversight jurisdictions, fraud facilitation.
    \item \textbf{NFT Layering:} Inflated and inconsistent pricing.
    \item \textbf{MSB Agents:} Operational and geographic risk exposure.
    \item \textbf{Over/Under-Invoicing:} Different pricing for identical goods.
    \item \textbf{Law Firm UBO:} Decline when UBO disclosure is refused.
    \item \textbf{Dual-Use Goods:} Suspended delivery and sanctions screening.
    \item \textbf{Law Firm STR:} File STR and delay for third-party payments.
    \item \textbf{FIU Guidance:} Assess exposure and update EWRA.
    \item \textbf{Guidance Impacts:} Customer experience, control redundancy, resources.
    \item \textbf{Guidance Effectiveness:} Consistent, aligned, demonstrable compliance.
    \item \textbf{Legal Sector Training:} Training and escrow risk assessment.
    \item \textbf{PEP Indirect Ownership:} Enhanced due diligence required.
\end{enumerate}
\end{block}

\vspace{0.05cm}
\tipbox{\footnotesize \textbf{FINAL LESSON:} Every topic teaches us that understanding the \textit{underlying principle} is more important than memorizing facts. Know \textit{why} and you will know \textit{what} to do.}
\end{frame}

% ================================================================
% SLIDE 26: THANK YOU
% ================================================================
\begin{frame}
\centering
\vspace{2.5cm}
{\fontsize{22}{26}\selectfont \textbf{Thank You}} \\[0.3cm]
{\fontsize{14}{17}\selectfont Keep learning. Keep growing.} \\[0.2cm]
\rule{3cm}{1pt} \\[0.2cm]
{\fontsize{10}{12}\selectfont \textit{Compliance is not about perfection — it is about continuous improvement.}}
\vspace{0.5cm}
{\fontsize{9}{11}\selectfont \textcolor{darkgray}{End of Tutorial}}
\end{frame}

\end{document}